% concatenated from DombosPartitionCongruences_daSilvaSellers_BSMM_15June2026.tex
\documentclass{amsart}

\usepackage{tikz}
\usepackage{xcolor}

\usepackage{latexsym,amsmath,amssymb,epic}
\usepackage{latexsym,amssymb,epic}
\usepackage{amsmath,amsthm}
\usepackage{color}


%%
%% JAS - I've inserted the package soul for now in order to allow for strikethrough to be done during the editing process... so that we can show one another where we suggest crossing out a word or phrase 
\usepackage{soul}
%%
%%

\usepackage{hyperref}
\hypersetup{
	colorlinks=true, %set true if you want colored links
	linktoc=all, %set to all if you want both sections and subsections linked
	linkcolor=blue} %choose some color if you want links to stand out

\newtheorem{theorem}{Theorem}[section]
\newtheorem{lemma}[theorem]{Lemma}
\newtheorem{corollary}[theorem]{Corollary}
\newtheorem{proposition}[theorem]{Proposition}
\newtheorem{definition}[theorem]{Definition}
\newtheorem{conjecture}[theorem]{Conjecture}
\newtheorem{example}[theorem]{Example}
\newtheorem{remark}[theorem]{Remark}

\begin{document}
	
	
	\title[A Partition Function of Dombos]{Arithmetic Properties Satisfied by a Recent Integer Partition Function of Dombos}
	
	
	\author{Robson da Silva}
	\address{Universidade Federal de S\~ao Paulo, S\~ao Jos\'e dos Campos, SP 12247--014, Brazil}
	\email{silva.robson@unifesp.br}

	\author{James A. Sellers}
	\address{Department of Mathematics and Statistics, University of Minnesota Duluth, Duluth, MN 55812, USA}
	\email{jsellers@d.umn.edu}

	
	
	\subjclass[2010]{11P83, 05A17}
	
	\keywords{partitions, grounded partitions, congruences, generating functions, $q$--series, dissections}
	
	%\footnotetext[1]{} \vspace{0.5in}
	%\footnotetext[1]{2010 AMS Classification Numbers: Primary, 11P81; Secondary, 05A17.}

\begin{abstract}
In recent work of Dombos, the set of integer partitions of $n$ wherein the parts are either divisible by 4 or congruent to $\pm 1 \pmod{6}$ arose in a natural way.   In this work, we will denote the function which counts the number of such partitions of $n$ by $dp(n)$.  Using elementary generating function manipulations and classical $q$--series results, we prove several congruences satisfied by $dp(n)$.  As an example, we prove that, for all $\alpha \geq 1$ and all $n \geq  0$,
\begin{equation*}
    dp \left( 3^{2\alpha + 1}n + \frac{7 \cdot 9^\alpha + 1}{4} \right) \equiv 0 \pmod{3}.
\end{equation*} 


\end{abstract}

	\maketitle
	
	%\noindent{\footnotesize{\bf Abstract.}
	


\section{Introduction} 

A partition of a positive integer $n$ is a nonincreasing sequence of positive integers $\lambda_1, \lambda_2, \dots, \lambda_r$ such that $n=\lambda_1+ \lambda_2+ \dots+ \lambda_r$.  We say that each $\lambda_i$, $1\leq i\leq r$, is called a part of the partition.  We denote the number of partitions of $n$ by $p(n)$.  

In the late 1910's, Ramanujan \cite{Ramanujan} proved the following congruences satisfied by $p(n)$ for the moduli $5, 7,$ and $11$ using generating function manipulations. 

\begin{theorem}%\label{ramanujanCongruences}
    For all $n\geq 0$, 
    \begin{align*}
        p(5n+4) &\equiv 0 \pmod{5}, \\
        p(7n+5) &\equiv 0 \pmod{7}, \text { and}\\
        p(11n+6) &\equiv 0 \pmod{11}.
    \end{align*}
\end{theorem}
These congruences, and others that have since been proven for numerous restricted partition functions, serve as the motivation for this work.  

In recent work of Dombos \cite{Dombos}, the set of integer partitions of $n$ wherein the parts are either divisible by 4 or congruent to $\pm 1 \pmod{6}$ arose in a natural way.  (See \cite[Theorem 1.6]{Dombos} for more details.)  In this work, we will denote the function which counts the number of such partitions of $n$ by $dp(n)$.  In light of this definition, it is clear that the generating function for $dp(n)$ is given by the following:  
\begin{equation}
G(q):=\sum_{n=0}^{\infty} dp(n)q^n = \frac{1}{(q;q^6)_\infty (q^5;q^6)_\infty (q^4;q^4)_\infty},
\label{gf1}
\end{equation}
where $(a;q)_\infty := (1-a)(1-aq)(1-aq^{2})\cdots$
is the usual $q$-Pochhammer symbol. 

The primary goal of this brief note is to prove a number of Ramanujan--like congruences satisfied by the function $dp(n)$ in certain arithmetic progressions.  In the work below, we will prove the following results. 

\begin{theorem}
For all $n \geq 0$,
	\begin{align}
	dp(6n+4) & \equiv 0 \pmod{2}, \label{c1} \\ 
%    dp(27n+16) & \equiv 0 \pmod{3}, \label{c2} \\
    dp(18n+10) & \equiv 0 \pmod{4}, \label{c3} \\
    dp(54n+52) &  \equiv 0 \pmod{8}. \label{c4}
    \end{align}
\label{Th1}
\end{theorem}

\begin{theorem} Let $p$ be a prime such that $p \equiv 17 \mbox{ or } 23 \pmod{24}$. Then for all $k,m \geq 0$ with $p \nmid m$, we have
\begin{equation*}
dp\left( 6p^{2k+1}m + \frac{15p^{2k+2}+1}{4} \right) \equiv 0 \pmod{4}.
\end{equation*}
\label{Th2}
\end{theorem}

\begin{theorem}
For all $n \geq 0$,
\begin{align}
dp(27n+16) & \equiv 0 \pmod{3}.
\label{c2}
\end{align}
\label{Th3}
\end{theorem}

\begin{theorem}
For all $n \geq 0$,
\begin{equation*}
dp(27n+7) \equiv dp(3n+1) \pmod{3}.
\end{equation*}
\label{internal_cong}
\end{theorem}

\begin{theorem}
For all $\alpha \geq 1$ and all $n \geq  0$,
\begin{equation*}
    dp \left( 3^{2\alpha + 1}n + \frac{7 \cdot 9^\alpha + 1}{4} \right) \equiv 0 \pmod{3}.
\end{equation*}
\label{Th4}
\end{theorem}

	







\section{Preliminaries}

Throughout the remainder of this paper, we define 
		$$f_k := (q^k;q^k)_{\infty}$$
in order to shorten the notation.

We begin by rewriting the generating function \eqref{gf1} in the following form.

\begin{lemma}
\begin{equation}
\sum_{n=0}^{\infty} dp(n)q^n = \frac{f_2f_3}{f_1f_4f_6}.
\label{gf2}
\end{equation}
\end{lemma}

\begin{proof}
Initially, we note that
\begin{align*}
\frac{1}{(q;q^6)_\infty (q^5;q^6)_\infty} & = \frac{(q^2;q^6)_\infty(q^3;q^6)_\infty(q^4;q^6)_\infty(q^6;q^6)_\infty}{(q;q^6)_\infty(q^2;q^6)_\infty(q^3;q^6)_\infty(q^4;q^6)_\infty(q^5;q^6)_\infty(q^6;q^6)_\infty} \\
& = \frac{(q^3;q^3)_\infty(q^2;q^6)_\infty(q^4;q^6)_\infty}{(q;q^3)_\infty(q^2;q^3)_\infty(q^3;q^3)_\infty} \\
& = \frac{(q^2;q^6)_\infty(q^4;q^6)_\infty}{(q;q^3)_\infty(q^2;q^3)_\infty} \\
& = \frac{(q^2;q^2)_\infty}{(q^6;q^6)_\infty}\cdot \frac{(q^3;q^3)_\infty}{(q;q)_\infty} \\
& = \frac{f_2f_3}{f_1f_6}.  
%& = \frac{(-q;q^3)_\infty(-q^2;q^3)_\infty(q;q^3)_\infty(q^2;q^3)_\infty}{(q;q^3)_\infty(q^2;q^3)_\infty} \\
%& = (-q;q^3)_\infty(-q^2;q^3)_\infty.
\end{align*}
Thus, it follows from \eqref{gf1} that
\begin{align*}
 \sum_{n=0}^{\infty} dp(n)q^n 
 %& = \frac{1}{f_4} (-q;q^3)_\infty(-q^2;q^3)_\infty \\
 %& = \frac{1}{f_4} \frac{(-q;q^3)_\infty(-q^2;q^3)_\infty (-q^3;q^3)_\infty}{(-q^3;q^3)_\infty} \\
 %& = \frac{1}{f_4} \frac{(-q;q)_\infty}{(-q^3;q^3)_\infty} \\
 & = \frac{f_2f_3}{f_1f_4f_6},
\end{align*}
which is \eqref{gf2}.
\end{proof}

The proofs of the main results mentioned above require a few well-known 2- and 3-dissections.

\begin{lemma}[\cite{R}, Lemma 1]
We have
\begin{align}
{f_{1}^{2}} & =  \frac{f_{2}f_8^5}{f_{4}^{2}f_{16}^{2}} -2q\frac{f_{2}f_{16}^{2}}{f_{8}}, \label{eq1061} \\
\frac{1}{f_{1}^{2}} & =  \frac{f_{8}^{5}}{f_{2}^{5}f_{16}^{2}} + 2q\frac{f_{4}^{2}f_{16}^{2}}{f_{2}^{5}f_{8}}.   
\label{eq2}
\end{align}
\end{lemma}

%\begin{lemma}[\cite{H}, Eq. (30.10.3)]
%We have
%\begin{equation}
%\frac{f_3}{f_1} = \frac{f_4f_6f_{16}f_{24}^2}{f_2^2f_8f_{12}f_{48}} + q \frac{f_6f_8^2f_{48}}{f_2^2f_{16}f_{24}}.   
%\label{eq3}
%\end{equation}
%\end{lemma}



%\begin{lemma}[\cite{Yao}, Eq. (2.19)]
%We have
%\begin{equation}
%\frac{f_{4}}{f_{1}}  = \displaystyle\frac{f_{12}f_{18}^{4}}{f_{3}^3f_{36}^2} + q\displaystyle\frac{f_{6}^2f_{9}^{3}f_{36}}{f_{3}^4f_{18}^2} + 2q^2 \displaystyle\frac{f_{6}f_{18}f_{36}}{f_{3}^3}.   
%\label{eq4}
%\end{equation}
%\end{lemma}

\begin{lemma}[\cite{Toh}, Eq. (2.1c)]
We have
\begin{equation}
\frac{f_{2}}{f_{1}f_{4}}  = \displaystyle\frac{f_{18}^{9}}{f_{3}^2f_{9}^3f_{12}^2f_{36}^3} + q\displaystyle\frac{f_{6}^2f_{18}^{3}}{f_{3}^3f_{12}^3} +q^2 \displaystyle\frac{f_{6}^4f_{9}^3f_{36}^{3}}{f_{3}^4f_{12}^4f_{18}^3}.   
\label{eq5}
\end{equation}
\end{lemma}

\begin{lemma}[\cite{Toh}, Eq. (2.1b)]
We have
\begin{equation}
\frac{f_{2}}{f_{1}^2} = \frac{f_{6}^4f_{9}^6}{f_{3}^8f_{18}^3} + 2q\displaystyle\frac{f_{6}^3f_{9}^3}{f_{3}^7} + 4q^2\displaystyle\frac{f_{6}^2f_{18}^3}{f_{3}^6}.   
\label{eq6}
\end{equation}
\end{lemma}

\begin{lemma}[\cite{B1}, Eq. (22.4)]
We have
\begin{equation}
\frac{f_{1}^2}{f_{2}} = \displaystyle\frac{f_{9}^2}{f_{18}} - 2q\displaystyle\frac{f_{3}f_{18}^2}{f_{6}f_{9}}.   
\label{eq7}
\end{equation}
\end{lemma}



%\begin{lemma}[\cite{H}, Eq. (14.8.5)]
%We have
%\begin{equation}
%f_{1}^{3}  =  \displaystyle\frac{f_{6}f_{9}^6}{f_{3}f_{18}^3} + 4q^3\displaystyle\frac{f_{3}^2f_{18}^{6}}{f_{6}^2f_{9}^3} - 3qf_{9}^{3}.   
%\label{eq8}
%\end{equation}
%\end{lemma}

\begin{lemma}[\cite{Toh}, Eq. (3.1)] We have
\begin{equation}
\displaystyle\frac{f_{2}^3}{f_{1}^3} = \displaystyle\frac{f_{6}}{f_{3}} +3q\displaystyle\frac{f_{6}^4f_{9}^{5}}{f_{3}^8f_{18}} +6q^2\displaystyle\frac{f_{6}^3f_{9}^{2}f_{18}^{2}}{f_{3}^7} +  12q^3\displaystyle\frac{f_{6}^2f_{18}^5}{f_{3}^6f_{9}}.  
\label{eq1062}
\end{equation}
\end{lemma}

We recall Ramanujan's theta functions
\begin{align*}
	f(a,b) & :=\sum_{n=-\infty}^\infty a^\frac{n(n+1)}{2}b^\frac{n(n-1)}{2}, \mbox{ for } |ab|<1, \nonumber \\
	\phi(q) & := f(q,q) = \sum_{n=-\infty}^{\infty} q^{n^2} = \frac{f_2^5}{f_1^2f_4^2}, \\
	\psi(q) & := f(q,q^3) = \sum_{n=0}^{\infty} q^{n(n+1)/2} = \frac{f_2^2}{f_1}. 
\end{align*}
With the above definitions, we can state the following 3-dissection of $1/\psi(q)$.

%\begin{lemma}[\cite{Shane}, Eq. (14)]
%We have
%\begin{align}
%\frac{1}{f_{1}^{3}} & = \displaystyle\frac{f_{9}^{3}}{f_{3}^{12}} \left( f_{3}^2 \frac{\phi(-q^9)^6}{\phi(-q^3)^2} + 8q^3f_{3}^2 \frac{\phi(-q^9)^3\psi(q^9)^3}{\phi(-q^3)\psi(q^3)} + 16q^6f_{3}^2\frac{\psi(q^9)^6}{\psi(q^3)^2} \right. \label{eq9} \\
%& \ \ \ \ \ \ \ \ \ \ \left. + 3qf_{3}f_9^3\frac{\phi(-q^9)^3}{\phi(-q^3)} + 12q^4f_{3}f_9^3 \frac{\psi(q^9)^3}{\psi(q^3)} +9q^2f_9^6 \right).  \nonumber 
%\end{align}
%\end{lemma}

\begin{lemma}[\cite{Shane}, Lemma 2.2] We have
	\begin{equation}
	\frac{1}{\psi(q)} = \frac{f_3^2f_9^3}{f_6^6} - q\frac{f_3^3f_{18}^3}{f_6^7} + q^2\frac{f_3^4f_{18}^6}{f_6^8f_9^3}.
	\label{3diss1/psi}
	\end{equation}
\end{lemma}

Lastly, we require the classical identity of Jacobi (see \cite[Theorem 1.3.9]{B1})
\begin{equation}
	f_1^3 = \sum_{n=0}^{\infty} (-1)^n (2n+1) q^{n(n+1)/2}.
\label{Jacobi}
\end{equation}
%and Euler's identity (see \cite[Eq. (1.6.1)]{H})
%\begin{equation}
%	f_1 = \sum_{n=-\infty}^{\infty} (-1)^n q^{n(3n-1)/2}.
%	\label{Euler}
%\end{equation}
%%%%%%%%%%%%%%%%%%%%%%%%%%%%%%%%%%%%
%%%%%%%%%%%%%%%%%%%%%%%%%%%%%%%%%%%%
%%%%%%%%%%%%%%%%%%%%%%%%%%%%%%%%%%%%


\section{Proof of Theorem \ref{Th1}}

Let 
$$H(q) = \sum_{n=0}^{\infty}h(n)q^n = G(-q) = \sum_{n=0}^{\infty}dp(n)(-q)^n.$$
Thus, $h(n) = (-1)^ndp(n)$ for all $n$. We note that
$$(-q;-q)_\infty = \frac{f_2^3}{f_1f_4}.$$
Using this identity and \eqref{gf2} we find that
\begin{equation*}
H(q) = \frac{f_1f_6^2}{f_2^2f_3f_{12}} = \frac{f_6^2}{f_3f_{12}} \frac{1}{\psi(q)}.
\end{equation*}
Thanks to \eqref{3diss1/psi} we obtain
\begin{align*}
\sum_{n=0}^{\infty}h(3n+1)q^{3n+1} & = -q \frac{f_3^2f_{18}^3}{f_6^5f_{12}},
\end{align*}
which, after dividing by $q$ and replacing $q^3$ by $q$, yields
\begin{align*}
\sum_{n=0}^{\infty}h(3n+1)q^{n} & = - \frac{f_1^2f_{6}^3}{f_2^5f_{4}}.
\end{align*}   
Using \eqref{eq1061}, we can 2-dissect the expression above to obtain
\begin{align*}
\sum_{n=0}^{\infty}h(6n+4)q^{2n+1} & = 2q \frac{f_6^3f_{16}^2}{f_2^4f_{4}f_8}.
\end{align*}
Thus,
\begin{align*}
\sum_{n=0}^{\infty}h(6n+4)q^{n} & = 2 \frac{f_3^3f_{8}^2}{f_1^4f_{2}f_4},
\end{align*}
which proves \eqref{c1} since $dp(6n+4) = h(6n+4) \equiv 0 \pmod{2}$.

Working modulo 8, we have
\begin{align}
\sum_{n=0}^{\infty}h(6n+4)q^{n} & = 2 \frac{f_3^3f_{8}^2}{f_1^4f_{2}f_4} \equiv 2 f_3^3\frac{f_4^3}{f_2^3} \pmod{8}.
\label{eq1063}
\end{align}
Thanks to \eqref{eq1062} we have
\begin{align*}
\sum_{n=0}^{\infty}h(18n+10)q^{3n+1} & \equiv 4q^4 \frac{f_3^3f_{12}^3f_{18}^2f_{36}^2}{f_6^7} \pmod{8},
\end{align*}
which simplifies to
\begin{align*}
\sum_{n=0}^{\infty}h(18n+10)q^{n} & \equiv 4q \frac{f_1^3f_{4}^3f_{6}^2f_{12}^2}{f_2^7} \pmod{8}.
\end{align*}
It follows that $h(18n+10) \equiv 0 \pmod{4}$, which completes the proof of \eqref{c3} since $h(n) = (-1)^n dp(n)$.

In order to prove \eqref{c4}, we use \eqref{eq1062} to extract the terms of the form $q^{3n+2}$ from \eqref{eq1063}. The resulting congruence is
\begin{align}
\sum_{n=0}^{\infty}h(18n+16)q^{n} & \equiv 6\frac{f_1^3f_{4}^4f_{6}^5}{f_2^8f_{12}}  \pmod{8}\nonumber \\
& \equiv 6\frac{f_6^5}{f_{12}} f_1^3 \pmod{8} \nonumber \\
& \equiv 6\frac{f_6^5}{f_{12}} \sum_{n=0}^{\infty} (-1)^n (2n+1) q^{n(n+1)/2} \pmod{8} \label{eq1064}
\end{align}
by \eqref{Jacobi}. We note that $f_6^5/f_{12}$ is a function of $q^3$ and $n(n+1)/2$ is never congruent to 2 modulo 3. Hence, once we expand \eqref{eq1064} as a power series in $q$, there will be no powers of the form $q^{3n+2}$. Therefore, for all $n \geq 0$, 
$$h(18(3n+2)+16) = h(54n+52) \equiv 0 \pmod{8},$$
which completes the proof of \eqref{c4}.



%%%%%%%%%%%%%%%%%%%%%%%%%%%%%%%%%%%%%%%%%%%
%%%%%%%%%%%%%%%%%%%%%%%%%%%%%%%%%%%%%%%%%%%
%%%%%%%%%%%%%%%%%%%%%%%%%%%%%%%%%%%%%%%%%%%


\section{Proof of Theorem \ref{Th2}}

Initially, we use \eqref{eq5} to extract the terms of the form $q^{3n+1}$ from \eqref{gf2}, which gives
$$\sum_{n=0}^{\infty} dp(3n+1)q^{3n+1} = q\frac{f_{6}f_{18}^3}{f_3^2f_{12}^3}.$$
After dividing by $q$ and replacing $q^3$ by $q$, we are left with
\begin{align}
\sum_{n=0}^{\infty} dp(3n+1)q^{n} & = \frac{f_{2}f_{6}^3}{f_1^2f_{4}^3}. \label{eq3n+1}
\end{align}

Now we use \eqref{eq2} to extract the odd parts of \eqref{eq3n+1}, which gives us
\begin{align}
\sum_{n=0}^{\infty} dp(6n+4)q^{n} & = 2\frac{f_3^3f_8^2}{f_1^4f_2f_4} \nonumber \\
& \equiv 2f_2^3 f_3^3\pmod{4}. \label{eq28.5.1}
\end{align}

Using \eqref{Jacobi}, it follows from \eqref{eq28.5.1} that
\begin{align*}
\sum_{n=0}^{\infty} dp(6n+4)q^{n} & \equiv 2f_2^3f_3^3 \pmod{4}\\
& = 2 \sum_{k,l=1}^{\infty} (-1)^{k+l}(2k+1)(2l+1)q^{k(k+1)+3l(l+1)/2} \\
& \equiv 2\sum_{k,l=1}^{\infty} q^{k(k+1)+3l(l+1)/2} \pmod{4}.
\end{align*}
Thus, 
\begin{align*}
\sum_{n=0}^{\infty} dp(6n+4)q^{24n+15} &  \equiv 2\sum_{k,l=1}^{\infty} q^{6(2k+1)^2+(6l+3)^2} \pmod{4}.
\end{align*}
It follows that $dp(6n+4) \equiv 0 \pmod{4}$ if $24n+15$ is not of the form $6x^2+y^2$. 

Since $p \equiv 17 \mbox{ or } 23 \pmod{24}$, it follows that $\displaystyle \left( \frac{-6}{p}\right) = -1$, which implies that $\nu_p(N)$ is even if $N$ is of the form $6x^2+y^2$. Taking $n = p^{2k+1}m + 5({p^{2k+2}-1})/{8}$, we have 
$$24n+15 = 24p^{2k+1}m + 15p^{2k+2} = p^{2k+1}(24m+15p).$$
Therefore, $\nu_p(24n+15)$ is odd and the result follows by the fact that
$$6\left( p^{2k+1}m + \frac{5(p^{2k+2}-1)}{8} \right) + 4 = 6p^{2k+1}m + \frac{15p^{2k+2}+1}{4}.$$

\section{Proofs of Theorems \ref{Th3}--\ref{Th4}}
We now transition to proving Theorems \ref{Th3}--\ref{Th4} where the focus in on arithmetic properties modulo 3.  Let us begin by proving \eqref{c2}. It follows from \eqref{eq3n+1} that
\begin{align*}
\sum_{n=0}^{\infty} dp(3n+1)q^{n} & \equiv \frac{f_{2}f_{6}^3}{f_1^2f_{12}} \pmod{3}. 
\end{align*}
Thanks to \eqref{eq6}, we get
\begin{align*}
\sum_{n=0}^{\infty} dp(9n+7)q^{3n+2} \equiv 4q^2\frac{f_{6}^5f_{18}^3}{f_3^6f_{12}} \pmod{3}.
\end{align*}
Dividing by $q$ and replacing $q^3$ by $q$ yields 
\begin{align*}
\sum_{n=0}^{\infty} dp(9n+7)q^{n} \equiv \frac{f_{2}^5f_{6}^3}{f_1^6f_{4}} \equiv \frac{f_6^4f_2^2}{f_3^2f_4} \pmod{3}.
\end{align*}
Finally, we use \eqref{eq7} to get
\begin{align}
\sum_{n=0}^{\infty} dp(9n+7)q^{n} \equiv \frac{f_6^4}{f_3^2} \left( \frac{f_{18}^2}{f_{36}} - 2q^2\displaystyle\frac{f_{6}f_{36}^2}{f_{12}f_{18}} \right) \pmod{3},
\label{eq29.5.1}
\end{align}
from which \eqref{c2} follows.

We next consider the ``internal congruence'' given in Theorem \ref{internal_cong}.  It follows from \eqref{eq29.5.1} that
\begin{align*}
\sum_{n=0}^{\infty} dp(27n+7)q^{3n} & \equiv \frac{f_6^4f_{18}^2}{f_3^2f_{36}}  \pmod{3}.
\end{align*}
Replacing $q^3$ by $q$, we are left with
\begin{align*}
\sum_{n=0}^{\infty} dp(27n+7)q^{n} & \equiv \frac{f_2^4f_{6}^2}{f_1^2f_{12}} \\
& \equiv \frac{f_2f_{6}^3}{f_1^2f_{4}^3} \pmod{3}.
\end{align*}
Thus, it follows by \eqref{eq3n+1} that
\begin{align*}
\sum_{n=0}^{\infty} dp(27n+7)q^{n}  \equiv \sum_{n=0}^{\infty} dp(3n+1)q^{n} \pmod{3},
\end{align*}
which completes the proof.

Lastly, we use the previous two theorems to prove Theorem \ref{Th4}.
We prove this theorem by induction on $\alpha$.  Note that the case $\alpha=1$ has already been established in Theorem \ref{Th3} above.  We now use it as our base case in a proof by induction.   
Thus, we assume that, for some $\alpha\geq 1$ and all $n\geq 0$, 
$$
dp \left( 3^{2\alpha + 1}n + \frac{7 \cdot 9^\alpha + 1}{4} \right) \equiv 0 \pmod{3}.
$$
We then want to prove that 
$$
dp \left( 3^{2\alpha + 3}n + \frac{7 \cdot 9^{\alpha+1} + 1}{4} \right) \equiv 0 \pmod{3}.
$$

We note that
\begin{align*}
dp \left( 3 \left( 3^{2\alpha}n + \frac{3\cdot 7 \cdot 9^{\alpha-1} -1}{4} \right) +1 \right) & = dp \left(  3^{2\alpha+1}n + \frac{7 \cdot 9^{\alpha} -3}{4} +1 \right) \\
& = dp \left(  3^{2\alpha+1}n + \frac{7 \cdot 9^{\alpha} +1}{4} \right) \\
& \equiv 0 \pmod{3},
\end{align*}
by the induction hypothesis. Finally, by Theorem \ref{internal_cong}, we have
\begin{align*}
dp \left( 3^{2\alpha + 3}n + \frac{7 \cdot 9^{\alpha+1} + 1}{4} \right) & = dp \left( 3^{2\alpha + 3}n + \frac{7 \cdot 9^{\alpha+1} -27}{4} + 7 \right) \\
& = dp \left( 27 \left( 3^{2\alpha}n + \frac{3 \cdot 7 \cdot 9^{\alpha-1} - 1}{4} \right) + 7 \right) \\
& \equiv dp \left( 3 \left( 3^{2\alpha}n + \frac{3\cdot 7 \cdot 9^{\alpha-1} -1}{4} \right) +1 \right) \pmod{3}\\
& \equiv 0 \pmod{3},
\end{align*}
which concludes the inductive proof.

\begin{thebibliography}{00}
\bibitem{B1} B. C. Berndt, {\it Ramanujan's notebooks}, Part III, Springer, 1991.
\bibitem{Shane} S. Chern and L.-J. Hao, {Congruences for partition functions related to mock theta functions}, Ramanujan J. {\bf 48} (2019), 369--384.
\bibitem{R} R. da Silva and J. A. Sellers, {\it New congruences for 3-regular partitions with designated summands}, Integers {\bf 20A} (2020), A6.
\bibitem{Dombos}  B. Dombos, {\it A new grounded partition identity of type $D_4^{(3)}$}, SIGMA {\bf 22} (2026), 042.  
%\bibitem{H} M. D. Hirschhorn, {The power of $q$, a personal journey}, Developments in Mathematics, v. 49, Springer, 2017.
\bibitem{Ramanujan} S. Ramanujan, {\it Some properties of $p(n)$, the number of partitions of $n$}, Proceedings of the Cambridge Philosophical Society, {\bf 19} (1919), 207--210.
%\bibitem{Sellers1} M. D. Hirschhorn and J. A. Sellers, {\it Arithmetic properties of partitions with odd parts distinct}. Ramanujan J. {\bf 22} (2010), 273--284.
\bibitem{Toh}  P.C. Toh, \textit{Ramanujan type identities and congruences for partition pairs}. Discrete Mathematics {\bf 312} (2012), 1244--1250.
%\bibitem{Yao} O.X.M. Yao, \textit{The relations between $N(a,b,c,d;n)$ and $t(a,b,c,d;n)$ and $(p,k)$-parametrization of theta functions}. J. Math. Anal. Appl. {\bf 453} (2017), 125--143.
\end{thebibliography}

\end{document}

